%
% Complete documentation on the extended LaTeX markup used for Insight
% documentation is available in ``Documenting Insight'', which is part
% of the standard documentation for Insight.  It may be found online
% at:
%
%     http://www.itk.org/

\documentclass{InsightArticle}


%%%%%%%%%%%%%%%%%%%%%%%%%%%%%%%%%%%%%%%%%%%%%%%%%%%%%%%%%%%%%%%%%%
%
%  hyperref should be the last package to be loaded.
%
%%%%%%%%%%%%%%%%%%%%%%%%%%%%%%%%%%%%%%%%%%%%%%%%%%%%%%%%%%%%%%%%%%
\usepackage[dvips,
bookmarks,
bookmarksopen,
backref,
colorlinks,linkcolor={blue},citecolor={blue},urlcolor={blue},
]{hyperref}
\usepackage{graphicx}
\usepackage{listings}
\usepackage{subfigure}
\title{InvertIntensityImageFilter}

%\release{1.00}

\author{Ga\"etan Lehmann}
\authoraddress{Unit\'e de Biologie du D\'eveloppement et de la Reproduction, Institut National de la Recherche Agronomique, 78350 Jouy-en-Josas, France}


\begin{document}
\lstset{language=c++}
\maketitle

\ifhtml
\chapter*{Front Matter\label{front}}
\fi

\begin{abstract}
\noindent
InvertIntensityImageFilter is a convenient filter to invert the intensity of an image.
\keyword{UnaryFunctorImageFilter}
\end{abstract}


\section{Description}
InvertIntensityImageFilter is a convenient filter to invert the intensity of an image. An image can already be inverted with ITK, but it requires to use a more complex filter (ShiftScaleImageFilter or IntensityWindowingImageFilter).

\section{Implementation}
InvertIntensityImageFilter is based on UnaryFunctorImageFilter and simply subtract the maximum value to all pixels before casting the result to the output pixel type.
\small \begin{verbatim}
TOutput result = static_cast<TOutput>( m_Maximum - x );
\end{verbatim} \normalsize
The maximum value is not computed by the filter to avoid iterating twice over the input image. It may be given by the user, but the default value (NumericTraits$<$ InputPixelType $>$::max()) must be correct in most of cases.
 
\section{Performances}
The performances of the different solutions have been measured on a linux box with an Athlon64 2800+ and repeated 100 fold. The inverted image is an 8 bits 3D image of size $371 \times 371 \times 34$. Results are reported in Table~\ref{perf}.

\begin{table}[htbp]
\centering
\begin{tabular}{cc}
\hline
filter & Execution time \\
\hline
\hline
InvertIntensityImageFilter & 0.0367 s \\
ShiftScaleImageFilter & 0.0842 s \\
IntensityWindowingImageFilter & 0.0876 s \\
\hline
\end{tabular}
\caption{Execution time.\label{perf}}
\end{table}


\section{Usage}
As usual, user has to include the header:
\small \begin{verbatim}
#include "itkInvertIntensityImageFilter.h"
\end{verbatim} \normalsize
and declare the type, instanciate the filter, and set the input image
\small \begin{verbatim}
typedef itk::InvertIntensityImageFilter< InputType, OuputType > InvertType;
InvertType::Pointer invert = InvertType::New();
invert->SetInput( filter->GetOutput() );
\end{verbatim} \normalsize
The maximum value can be set with the SetMaximum() method, and read with GetMaximum().
\small \begin{verbatim}
invert->SetMaximum( 80 );
\end{verbatim} \normalsize
The default is itk::NumericTraits$<$ PType $>$::max().

\section{Example}
This filter can be use to invert a binary image, a distance map, etc.

\begin{figure}[htbp]
\centering
\includegraphics[width=0.5\textwidth]{cthead1}
\caption{The input image.}
\end{figure}

\begin{figure}[htbp]
\centering
\includegraphics[width=0.5\textwidth]{test}
\caption{The inverted image.}
\end{figure}


\section{Acknowledgments}
We thank Dr Pierre Adenot and MIMA2 confocal facilities (\url{http://mima2.jouy.inra.fr})
for providing image samples.

\bibliographystyle{plain}
\bibliography{InsightJournal}
\nocite{ITKSoftwareGuide}

\end{document}
